\exnumber{09}
\extitle{UI III - Ex 10 - Micro-interactions}
\exfiles{ex09/Figma file, named Name\_FirstName\_DeliverableName\_Date\_VersionNumber}
\exauthorize{Those named in the exercise}
\makeheaderfiles

\textbf{Instructions:}

Browse the internet and find at least 5 different examples of feedback micro-interactions or micro-animations on websites or apps.

Be careful! Feedback is helpful to provide information to the user, not just showing a catchy design. You need to make the difference.

\textbf{Deliverable:}

\begin{itemize}
\item A Figma file gathering 5 examples of \textbf{feedback} micro-interactions or animations. In each example, show the different versions of the element.
\item Upload the Figma file on Github, labeled as follows: "Name\_FirstName\_DeliverableName\_Date\_VersionNumber". For example, it can be labeled as "Doe\_John\_MidFiPrototype\_05242024\_V1".
\item Add a title within the file.
\item Pay attention to the global aspect of the file. It must be clear, simple and easy to read. You can get inspiration on canva.com but don't overload your design with too much details!
\end{itemize}

\textbf{Resources:}

\begin{itemize}
\item Why use micro-interactions in your design \url{https://hike.one/update/why-use-micro-animations-in-your-design}
\item Micro-interactions: why, when,  and how to use them to boost the UX \url{https://uxdesign.cc/micro-interactions-why-when-and-how-to-use-them-to-boost-the-ux-17094b3baaa0}
\end{itemize}

\textbf{Tips:}

\begin{itemize}
\item Feedback is the situation when the app specifically sends a piece of information to the user, like: 
\item Payment accepted
\item Refreshed feed
\item Information saved
\item etc.
\end{itemize}

